%% ma: L12-B01-0027
%% lop: 12
%% chuong: 1
%% bai: 1
%% dang: 12.1.18
%% loai: TN4
%% muc: TH
%% dap_an: "B"
%% nguon: "(NBV)-ĐỀ SỐ 16-MỖI NGÀY 1 ĐỀ THI 2025.pdf (D16, MỖI NGÀY 1 ĐỀ - Đề số 16), Phần I Câu 5"
%% kien_thuc: "Bài 1 (cực trị của hàm số; đạo hàm hàm phân thức)"
%% trang_thai: cho-duyet
%% ly_do: "Dạng mới đề xuất 12.1.18 chờ giáo viên duyệt"
%% ghi_chu: "Đáp án gốc B đúng."
\begin{ex}
Trong các hàm số sau, hàm số nào không có cực trị?
\choice{$y=\dfrac{x^2}{x+1}$}{\True $y=\dfrac{x+2}{x+1}$}{$y=-x^3+3x+2$}{$y=x^2+2x$}
\loigiai{B: $y'=\dfrac{-1}{(x+1)^2}<0$ với mọi $x\ne-1$ nên hàm số không có cực trị.

A: $y'=\dfrac{x^2+2x}{(x+1)^2}$ đổi dấu tại $x=-2$ và $x=0$; C: $y'=-3x^2+3$ đổi dấu tại $x=\pm1$; D: $y'=2x+2$ đổi dấu tại $x=-1$. Các hàm số A, C, D đều có cực trị. Chọn B.}
\end{ex}
